% Build the ACM target with: make acm
% Build the portable fallback with: make plain
\newif\ifplainarticle
\ifdefined\PLAINARTICLE
  \plainarticletrue
\else
  \plainarticlefalse
\fi

% CPP 2027 uses lightweight double-blind review: build the anonymized
% submission with `make anon` (acmart [anonymous,review]); the arXiv /
% camera-ready build keeps the author block and the CRediT statement.
\newif\ifanonsubmission
\ifdefined\ANONSUBMISSION
  \anonsubmissiontrue
\else
  \anonsubmissionfalse
\fi

\ifplainarticle
  \documentclass[10pt,twocolumn]{article}
  \usepackage[a4paper,margin=0.72in,columnsep=0.24in]{geometry}
  \usepackage[numbers,sort&compress]{natbib}
\else
  \ifanonsubmission
    \documentclass[sigplan,10pt,anonymous,review]{acmart}
    \settopmatter{printfolios=true,printccs=false,printacmref=false}
  \else
    \documentclass[sigplan,screen]{acmart}
    \settopmatter{printacmref=false,printfolios=true}
  \fi
  \setcopyright{none}
  \renewcommand\footnotetextcopyrightpermission[1]{}
  \citestyle{acmnumeric}
\fi

\ifplainarticle
  % acmart supplies \Description for figure accessibility; the portable
  % fallback class does not, so absorb it there.
  \providecommand{\Description}[1]{}
\fi

\usepackage[T1]{fontenc}
\usepackage[utf8]{inputenc}
\usepackage{booktabs}
\usepackage{tabularx}
\usepackage{array}
\usepackage{xspace}
\ifplainarticle
  \usepackage{amsmath,amssymb}
\else
  \usepackage{amsmath}
\fi
\ifplainarticle
  \usepackage[hidelinks]{hyperref}
\fi

\newcommand{\solver}{\textsc{EQT2-S2}\xspace}
\newcommand{\lean}{Lean~4\xspace}
\newcommand{\op}{\mathbin{\diamond}}
\newcommand{\Fin}{\mathop{\mathrm{Fin}}}
\newcommand{\Nat}{\mathbb{N}}
\newcommand{\code}[1]{\texttt{#1}}
\newcommand{\prov}[2]{\begingroup\scriptsize\raggedright\nolinkurl{#1}\\[-1pt]\texttt{#2}\par\endgroup}
\newcolumntype{Y}{>{\raggedright\arraybackslash}X}
\newcolumntype{R}{>{\raggedleft\arraybackslash}p{0.65in}}

\newcommand{\paperabstract}{%
Proof-producing automation concentrates on establishing theorems, while
counterexample generation is usually the province of separate model-finding
tools.  Equational implication demands both polarities at once: every input
requires either a universal \lean proof or an explicit \lean countermodel, and
an unverified failure on one side must never be treated as evidence for the
other.  We present a dual-polarity certifying architecture in which ordered
superposition, structured algebraic search, bounded finite-model
construction, and infinite-carrier witness generation all compile into
first-class \lean terms of opposite logical shape, each checked against the
corresponding judge-generated \code{Goal} under a shared dependency policy.
The unit of success is a typed \lean certificate rather than a verdict:
search state, heuristics, and unverified labels are compiled away before a
result is submitted for acceptance, so defects in the untrusted search can
reduce coverage but cannot by themselves carry an invalid certificate across
the checking boundary.  Ordered superposition is one producer inside this
architecture.

We evaluate a frozen submission on 2,889 problems, obtaining 1,429 positive
\lean proofs and 1,460 negative \lean witnesses, every one accepted by the
judge.  The unchanged submission subsequently produced accepted certificates
for all 200 problems of an official set released after the freeze, a temporal
holdout.  Every reported certificate came from the local, non-stochastic
cascade, with no language-model and no external-prover calls.  The emitted
certificates remain compact and inexpensive to check, which keeps production
at this scale compatible with the deployment envelope.  A harder research-tier
corpus remains unsolved, delimiting the present coverage.  The paper makes no
completeness or cross-system superiority claim, and binds every empirical
table to immutable ledgers.}

\title{A Dual-Polarity Certifying Cascade for Equational Implication:\\
Lean Proofs and Countermodels under Resource Bounds}

\ifplainarticle\else
  \acmConference[CPP 2027]{Certified Programs and Proofs}{January 11--12, 2027}{Mexico City, Mexico}
  \acmYear{2027}
  \setlength{\headheight}{18.5pt}
\fi

% TODO(author order): author order is TBD.
% Author order agreed 2026-08-24 (order reflects contribution; CRediT below).
\ifplainarticle
  \ifanonsubmission
    \author{Anonymous submission}
  \else
    \author{Haobo Ma\\ChronoAI Pte.\ Ltd., Singapore
      \and Wenlin Zhang\\National University of Singapore
      \and Manuel Israel C\'azares\\Bytepro AI, Mazatl\'an, Mexico}
  \fi
  \date{}
\else
  \author{Haobo Ma}
  \affiliation{\institution{ChronoAI Pte.\ Ltd.}\country{Singapore}}
  \author{Wenlin Zhang}
  \affiliation{\institution{National University of Singapore}\country{Singapore}}
  \author{Manuel Israel C\'azares}
  \affiliation{\institution{Bytepro AI}\city{Mazatl\'an}\country{Mexico}}
\fi

\begin{document}
\ifplainarticle
  \maketitle
\begin{abstract}
\paperabstract
\end{abstract}
\else
\begin{abstract}
\paperabstract
\end{abstract}
  \maketitle
\fi

\ifplainarticle\else
\keywords{equational reasoning, superposition, finite model finding, Lean,
proof certificates, reproducibility}
\fi

\section{Introduction}

An equational theory of magmas has one binary operation and no assumed
associativity, identity, or cancellation law.  Given equations $E_1$ and
$E_2$, the implication problem asks whether every magma satisfying $E_1$ also
satisfies $E_2$.  Proof-producing provers principally target the positive
direction: they establish a theorem or reconstruct a proof found by an ATP.
Counterexample generation is usually exposed through a separate model finder.
Neither interface alone meets the implication problem's two-sided contract.
Each input requires one of two mathematically opposed objects, and failure to
find or validate one object is never evidence for the other polarity.

The SAIR Mathematics Distillation Challenge, organized by Damek Davis and
Terence Tao with the SAIR Foundation, makes that gap operational
\cite{competition,etp}.  A true answer must contain a \lean proof of the
universal implication; a false answer must exhibit a magma in which the
hypothesis holds and the conclusion fails.  The same judge checks both forms.
A plausible label, proof sketch, failed proof search, or failed model search
receives no credit~\cite{competition,competitionRules}.

We address this gap with a Lean-facing dual-polarity certificate architecture.
Ordered superposition compiles to a positive equality term; finite and
structured models compile to reflected existential witnesses; and infinite
Austin models compile to explicit mathematical theorems and a separating
tuple.  All three paths target the judge-generated \code{Goal} and the same
dependency policy.  The unit of success is thus a typed \lean artifact, not a
verdict.  Search state is compiled away, and judge acceptance is the only
commit point.

The producers run as a resource-bounded cascade in one Python file, under both
the per-problem Solo protocol and the globally budgeted Marathon protocol.
Across the principal public, temporal-holdout, and distribution-drill
evaluation, the frozen artifact produced 2,889 accepted artifacts: 1,429
positive proofs and 1,460 negative witnesses.  Its bytes were frozen before
the official stress set was released; the same bytes later produced 200/200
accepted certificates on that set.  All reported accepted results are
produced deterministically, with zero language-model and zero external-prover
calls.  The separate research-tier experiment in
Section~\ref{sec:holdout} records the boundary of this coverage.

This paper makes four contributions.
\begin{itemize}
  \item A Lean-facing dual-polarity architecture makes positive proofs and
  negative witnesses first-class terms under one judge-generated goal and
  dependency policy, with acceptance rather than a search verdict as commit.
  \item A certificate-preserving compiler translates superposition traces into
  compact core-Lean proofs.  It slices the reachable recipe DAG, reconstructs
  substitutions and one-hole contexts, lifts equalities with \code{congrArg},
  composes symmetry and transitivity, and emits shared ancestors once in
  topological order.
  \item Three negative compilers emit reflected finite tables, algebraically
  structured operations, and explicit infinite-carrier witnesses.  The infinite path carries the full mathematical
  argument because exhaustive checking over $\Nat$ is unavailable.
  \item Under strict resource bounds, the implementation produced thousands
  of judge-accepted artifacts, including a temporal holdout produced by the
  frozen bytes; the evidence separates this result from completeness or
  cross-system superiority claims.
\end{itemize}

This architecture occupies a different point from existing positive-proof
automation.  Duper supplies proof-producing superposition inside Lean, while
Lean-Auto, Lean on Vampire, and Isabelle reconstruction connect mature provers
to proof assistants~\cite{duper,leanAuto,leanVampire,slam}.  Those systems establish the
positive-proof baseline.  The architecture studied here is complementary: a
narrower equational engine becomes one producer in a two-polarity system that
also generates checked finite and infinite models under a single-file
deployment contract.  Section~\ref{sec:rw} gives the full
comparison.

\begin{table*}[t]
\caption{Positioning against the closest automation architectures.  ``Negative''
means an explicit countermodel certificate, not failure to prove.}
\label{tab:systems}
\small
\begin{tabularx}{\textwidth}{@{}lYYYY@{}}
\toprule
System & Search locus & Checked positive result & Negative certificate & Deployment distinction \\
\midrule
Duper~\cite{duper} & Lean tactic, dependent type theory & Lean proof & no & general proof-assistant integration \\
Lean-Auto~\cite{leanAuto} & external E/Vampire/Z3 & reconstructed Lean proof & no & multi-ATP interface \\
Lean on Vampire~\cite{leanVampire} & external Vampire & reconstructed Lean proof & no & imports mature ATP proofs \\
\textsc{slam}~\cite{slam} & Isabelle $\lambda$-superposition & Isabelle proof & no & higher-order / Sledgehammer backend \\
\solver & single Python file, unit equations & positive Lean equality proof & finite or infinite Lean witness & dual polarity, no external prover, hash-bound holdout \\
\bottomrule
\end{tabularx}
\end{table*}

Section~\ref{sec:certarch} defines the Lean certificate architecture before
Section~\ref{sec:arch} presents the Python search cascade that supplies it.
Sections~\ref{sec:false} and~\ref{sec:true} describe the countermodel and proof
searches; Sections~\ref{sec:hosted} and~\ref{sec:evaluation} cover deployment
and evaluation.  We then state the reproducibility boundary and position the
system relative to equational-theory, superposition, and proof-assistant work.

\section{Lean Certificate Architecture}
\label{sec:certarch}

The architectural interface is a pair of Lean types, not a pair of Boolean
labels.  For each parsed input the judge generates a verdict-specific
\code{Goal}.  Suppressing the bound variables inside the identities, the two
shapes are
\[
  \mathit{Goal}^{+} :=
  \forall (G : \mathsf{Type})\,[\mathsf{Magma}\;G],\ E_1(G)\to E_2(G)
\]
and
\[
  \mathit{Goal}^{-} :=
  \exists (G : \mathsf{Type})\,(_ : \mathsf{Magma}\;G),\
  E_1(G)\mathbin{\land}\neg E_2(G).
\]
The solver submits a definition named \code{submission : Goal}; it never
reconstructs a look-alike proposition from its own parser.  Consequently,
Lean elaboration checks the actual binder order, operation instance, and
polarity chosen by the judge.

\begin{figure*}[t]
\centering
\small
\setlength{\tabcolsep}{4pt}
\begin{tabular}{@{}c@{\ }c@{\ }c@{\ }c@{\ }c@{}}
\fbox{superposition recipe DAG} & $\longrightarrow$ &
\fbox{positive equality compiler} & $\longrightarrow$ &
$\forall G,\ E_1\to E_2$ \\
\fbox{finite / structured model} & $\longrightarrow$ &
\fbox{reflected witness compiler} & $\longrightarrow$ &
$\exists G,\ E_1\land\neg E_2$ \\
\fbox{infinite Austin model} & $\longrightarrow$ &
\fbox{explicit theorem compiler} & $\longrightarrow$ &
$\exists G,\ E_1\land\neg E_2$ \\
&&&& $\downarrow$ \\
\multicolumn{5}{c}{\fbox{elaboration $+$ kernel checking $+$ dependency policy}}
\end{tabular}
\caption{Three certificate-producing pipelines share one Lean acceptance
boundary.  Search-specific representations disappear before acceptance.}
\Description{A superposition recipe DAG, finite or structured model, and
infinite Austin model pass through three compilers to positive or negative
Lean goals, then converge at elaboration, kernel checking, and dependency
inspection.}
\label{fig:certificate-compilation}
\end{figure*}

\subsection{Positive equality compiler}

The positive producer does not export its internal refutation.  Search treats
the negated goal $u\neq v$ as a negative unit clause and closes it when the two
sides unify.  Compilation then walks backward from that closure and retains
only recipes reachable from the two final equality chains.  The resulting DAG
excludes irrelevant saturation state; a topological traversal emits every
shared ancestor exactly once.

Each recipe records its source and target equations, directions, overlap
position, and substitution.  The compiler reconstructs the instantiated terms
and the one-hole target context.  It turns the source equality into an equality
at the overlap with \code{congrArg}, selects symmetry when a recorded direction
requires it, and joins the lifted source equality to the target equality by
transitivity.  Demodulation uses the same construction with one-way matching.
Thus a search over a negative clause becomes a positive equality chain in
core Lean, expressed as locally typed \code{have} bindings followed by an
\code{exact} term.

Compilation balances source size against transparent replay.  Printing every
normalization micro-step repeats contexts and substitutions, whereas replacing
the whole derivation by a large tactic call reintroduces tactic behavior and
library imports.  The compiler therefore emits one named equality per selected
inference, with normalization represented by explicit transitivity as needed.
If the direct printer is rejected, robust re-emission wraps individual steps in
a small set of core closing alternatives and resubmits the same derivation; it
does not rerun superposition or change the theorem.  If search ends without a
direct closure, a separate lemma-pool path may emit a bounded selection of the
smallest proved consequences, but the resulting tactic must still close the
exact judge goal or the row remains unanswered.

For example, from $x\op y=x$ the goal $(a\op b)\op c=a$ is compiled by putting
the instance $a\op b=a$ in the context $q\mapsto q\op c$ and composing with
$a\op c=a$:
\begin{quote}
\small\ttfamily
import JudgeProblem\\
def submission : Goal := by\\
\ \ intro G \_ h a b c\\
\ \ have e0 : (a $\diamond$ b) $\diamond$ c = a $\diamond$ c :=\\
\ \ \ \ congrArg (fun q => q $\diamond$ c) (h a b)\\
\ \ exact e0.trans (h a c)
\end{quote}
Term ordering, unification code, and the refutation data structure do not occur
in this artifact.

\subsection{Negative witness compilers}

A finite or structured search result is compiled into the carrier, a total
$\mathsf{Magma}$ operation, a proof that every assignment satisfies $E_1$, and a
concrete assignment falsifying $E_2$.  For carriers through ten,
\code{finOpTable} embeds a JSON-like Cayley table and \code{decideFin!}
reflects the finite checks.  Its parser extracts individual decimal
characters, so it cannot faithfully encode multi-digit entries.  Larger
finite carriers and algebraically structured families instead receive a
transparent operation such as \code{submission.op}, defined by core arithmetic
or a packed-table lookup.  The submitted theorem refers to that operation
directly, allowing Lean to evaluate the same structure without the table
parser.

Infinite witnesses require a different compiler.  For a recognized Austin
pair it defines the selected operation on $\Nat$ (or the opposite operation),
instantiates a fixed sequence of elementary lemmas proving $E_1$ universally,
and supplies a small tuple proving $\neg E_2$.  No exhaustive tactic can decide
the universal premise over $\Nat$; the full mathematical argument is therefore
part of the certificate rather than a Python-side claim.

\subsection{Dependency policy and acceptance}

The permitted imports are the judge's problem module and, for reflected finite
witnesses, its finite-decision support modules; positive proofs and current
infinite witnesses otherwise use core definitions and equality eliminators.
The emitted source contains no admitted terms.  The dependency inspector
checks the transitive declaration and axiom closure of \code{submission}
against the judge policy.  Submission-local helpers are the documented way to
package operations and lemmas: they are not global dependencies, but their
transparent bodies are still elaborated, compiled, and traversed by the kernel.
In particular, the allow-list inspects direct constants in the submitted term;
the documented exemption for names in the \code{submission} namespace permits
local auxiliary definitions but does not exempt their bodies from kernel
checking.  The policy controls dependencies, not the truth of definitions.

Acceptance is a commit operation.  Python-side evaluation can discard a bad
table or malformed term cheaply, but passing such a check records only a
candidate.  The solver commits a row only after the official judge has
elaborated \code{submission : Goal}, kernel-checked it, and accepted its
dependency closure.  A rejected term has no evidential effect on either
polarity; the next producer may run, or the row remains unanswered.

The positive compiler maintains a more precise replay invariant: for every
retained equation $a=b$, the recipe DAG emits a Lean term whose type is exactly
$a=b$ under instances of the input hypothesis.  The base recipe is an input
hypothesis instance.  The overlap and demodulation cases follow from the
recorded substitution, one-hole context, \code{congrArg}, symmetry, and
transitivity.  Dependencies precede users in the acyclic emission order, so induction over
that order justifies the invariant; this metatheoretic argument is not itself
mechanized.  At the closure roots the two
chains meet at the unified negative goal, yielding the requested positive
equality.  This argument specifies the compiler; the judge still checks each
emitted instance and is the only trusted checker.

\begin{table*}[t]
\caption{Lean constraints, the certificate design each one forced, and the
observed consequence.  Every consequence restates a measurement bound
elsewhere in this paper; the table introduces none of its own.}
\label{tab:lean-design}
\scriptsize
\setlength{\tabcolsep}{4pt}
\begin{tabularx}{\textwidth}{@{}p{0.20\textwidth}p{0.32\textwidth}Y@{}}
\toprule
Challenge & Certificate response & Observed consequence \\
\midrule
Saturation traces are too large & Reachable-recipe slicing, topological
emission, and one binding per shared ancestor & Public and holdout certificate
size distributions in Table~\ref{tab:certcost} \\
Broad tactic imports are too expensive & Core-only positive proof with
explicit lemmas & Austin compilation fell from about 330 s to within 3.1 s \\
\code{finOpTable} cannot encode multi-digit entries & Transparent namespaced
operation over the larger carrier & The official judge accepted the recorded
larger-carrier witnesses \\
The searcher or printer may be wrong & Compile search state away; require a
term of the exact judge type & Kernel-rejected frontier candidates remained
unaccepted \\
The hosted toolchain changes & Keep certificates dependency-light and locally
typed & Full frozen-artifact revalidation under Lean 4.33.1 \\
\bottomrule
\end{tabularx}
\end{table*}

\section{Single-File Solver Architecture}
\label{sec:arch}

\subsection{Deployment contract}

The submission contract requires one \code{solver.py} file no larger than the
competition limit.  Solver v2.5 is 189,504 bytes, and therefore remains below
that limit while containing all search code and fixed data.  In Solo mode the
proxy sends one JSON problem and its budget to a fresh process and receives
service requests over line-delimited JSON.  In Marathon mode the runner supplies
a JSONL manifest and an append-only output path.  The same source supports both
modes, selected by the runner-owned Marathon environment
variable~\cite{competitionRules}.  Section~\ref{sec:certarch} defines the
shared typed interface and explains why Python search remains outside the
trusted boundary.

At both protocol entry points we normalize the input operation symbol.  The
official problem format permits either an asterisk or a diamond, but the
solver's term parser has one internal token.  Mapping the asterisk to
$\op$ before parsing preserves the semantics and is operationally necessary;
Section~\ref{sec:encoding} records the failure that revealed it.

\subsection{Cheapest-first cascade}

The Solo cascade is ordered by expected cost and by which branch benefits from
remaining time.

\paragraph{Stage A: shallow countermodels.}
The solver first tries tiny brute-force tables, scalar linear and affine
operations, vector-linear families, low-degree polynomial operations, and
recognized infinite-carrier Austin models~\cite{austin}.  The expanded structured families
add about 0.8 seconds on a true input in the recorded development measurement.
No irregular finite-model search is performed yet.

\paragraph{Stage B: direct true certificates.}
Syntactic singleton collapse and substitution-instance checks produce direct
proofs.  A specialized singleton-forced prover then receives an eight-second
cap, although applicable cases normally close during its initial search.  The
general superposition engine follows with a quick budget
\(
\min(6,\max(1,T/600))
\)
seconds for a process allowance $T$; this is six seconds under the recorded
Solo allowance.  Most provable public instances terminate within this cap.

\paragraph{Stage C: irregular models.}
A finite-model search over carriers four through ten receives eight seconds.
It propagates the universal instances of the hypothesis and Latin-square
constraints rather than enumerating complete tables.

\paragraph{Stage D: deep false search.}
Only residuals reach heavier polynomial and vector families, capped at sixty
seconds, followed by another finite-model pass capped at one hundred twenty
seconds.  Deep false search precedes the deepest proof pass so that a difficult
false row does not consume most of its allowance in proof saturation.

\paragraph{Stage E: proof search.}
The ordered-superposition engine follows the completion tradition of
equational provers~\cite{knuthBendix,bachmairGanzinger} and receives
\(
\min(600,\max(20,T/6))
\)
seconds, which is six hundred seconds under the Solo allowance.  The engine
repeatedly raises term-size and variable caps.  If no direct refutation is returned, a
small pool of derived equalities can be emitted as exact lemmas followed by a
kernel-checked closing tactic.  Solo also contains a last-resort language-model
fallback, but it was never called for any archived accepted row; its measured
history is reported in Section~\ref{sec:evaluation}.

\begin{figure*}[t]
\centering
\small
\setlength{\tabcolsep}{5pt}
\begin{tabular}{@{}c@{\ }c@{\ }c@{\ }c@{\ }c@{}}
\fbox{A: structured models} & $\rightarrow$ &
\fbox{B: direct proofs + quick G3} & $\rightarrow$ &
\fbox{C: finite-model search} \\
&& $\downarrow$ && \\
\fbox{accepted certificate} & $\leftarrow$ &
\fbox{official \lean judge} & $\leftarrow$ &
\fbox{D: deep models $\rightarrow$ E: anytime G3}
\end{tabular}
\vspace{3pt}

\begin{minipage}{0.94\textwidth}
\footnotesize
Each search box is untrusted and may return no candidate.  Every Solo arrow to
the left carries a complete positive or negative certificate; only the judge's
accepted response terminates the cascade.  Marathon executes the same boxes in
two fairness passes and defers judging until the append-only answer file is
closed.
\end{minipage}
\caption{Cascade and trust boundary.  The distinctive interface is
dual-polarity: proof search and model search share the same fail-closed commit
point.}
\Description{A five-stage cascade sends proof and countermodel certificates to
the official Lean judge; only judge acceptance commits an answer.}
\label{fig:cascade}
\end{figure*}

\begin{figure}[t]
\footnotesize\ttfamily
normalize both input equations\par
for stage in [A, B, C, D, E]:\par
\hspace*{1em}candidate := run(stage, remaining time)\par
\hspace*{1em}if candidate = none: continue\par
\hspace*{1em}reply := judge(candidate.polarity, candidate.code)\par
\hspace*{1em}if reply = accepted: return\par
return unanswered
\caption{Solo scheduling pseudocode.  Search never exports a cached Boolean
verdict; a rejected candidate has no effect on the opposite polarity.}
\Description{Pseudocode normalizes the equations, executes five stages in
order, submits each complete candidate to the judge, and stops only on
acceptance.}
\label{fig:scheduler}
\end{figure}

Every stage independently verifies its candidate before the solver stops.  A
rejected certificate proceeds to the next strategy; a successful search
without a valid emitted term is not counted as a solution.  This control flow
distinguishes search coverage from certificate robustness.

\subsection{Representation and dispatch}

Both branches share a small non-associative term representation.  Parsing
produces a binary tree with variables at leaves; no reassociation or
commutative normalization is performed.  Alpha-renaming maps variables to a
canonical local order for cache keys and family matching, but emitted proofs
retain a map back to the judge's binder order.  Operation count, variable
support, depth, and tree size are computed once and reused by scheduling,
countermodel probes, and proof search.

Dispatch relies on positive recognizers.  A direct proof stage
returns only after constructing a substitution or equality chain.  An Austin
or central-groupoid stage returns only when the hypothesis matches its stored
schema and a concrete evaluation separates the goal.  A family that is merely
likely from equation identifiers is not selected: identifiers are useful for
diagnostics, but the solver matches parsed laws.  This property supports
independent replay and handles inputs whose labels or variable names differ
from those in the public files.

Candidate validation is stratified.  Low-cost Python evaluation removes
malformed tables, wrong witnesses, and printer bugs before a judge call.  Judge feedback
then closes the semantic and dependency obligations.  The development harness
records both outcomes: a ``found'' countermodel rejected by \lean indicates an
emitter defect or an encoding mismatch, whereas the absence of a Python
witness within a deadline indicates a coverage limitation.  This distinction informed the large-carrier encoding correction in
Section~\ref{sec:allowlist}.

\subsection{Solo and Marathon}

Solo gives each row a fresh process and a fixed allowance.  Marathon instead
requires triage across a whole manifest.  The batch path stably sorts problems
by a structural cost derived from operation occurrences, variable count, and
equation length.  It then runs two passes.  The first gives every row the
shallow stages before any residual receives deep work.  The second revisits
only residuals.  A fair cap is recomputed from remaining wall-clock time
divided by remaining rows and clamped to the documented per-row interval.  At most fifty-five
percent of a residual's current slice goes to the false branch; the proof
branch receives the rest.

Marathon outputs constitute durable state.  Each answer is one JSON line, flushed and
synchronized before the next row.  If the runner terminates the process, all
complete earlier lines remain scoreable.  A process-level timer bounds legacy
search loops that predate the batch driver.  The batch path leaves a residual
unanswered when its certificate-producing stages are exhausted.

\subsection{Cascade invariants and failure isolation}

Stage ordering does not assert that either polarity is more likely.  Instead,
it minimizes expected certificate cost while reserving time for both
branches.  A structured false witness is attractive early because checking a
closed operation is usually cheaper than saturation.  Direct true patterns are
equally attractive because they avoid both table construction and general proof
search.  Measurements showed that most provable rows close near the start of
saturation, which motivates placing the quick G3 pass before irregular model
search.  Deep false search precedes deep G3 only after every low-cost recognizer
has failed.

The implementation maintains a cascade invariant: before entering a stage,
the solver has no accepted certificate, and every earlier candidate has either
been absent, failed local construction, or been rejected by the judge.  It does
not maintain the stronger and generally invalid invariant that earlier stages
have ruled out their polarity.  Failure to find a finite model does not support
truth; failure to derive the goal does not support falsity.  Each stage can
therefore be analyzed and tested as a partial certificate producer, and the
composition is sound whenever the judge is sound, regardless of the producers'
coverage.

Exceptions are contained at strategy boundaries in the Marathon path so that a
malformed residual cannot discard previously completed rows.  Solo has a fresh
process per row and can rely on the runner to record a crash.  Within a
strategy, deadlines are checked at natural allocation points: carrier changes,
model-search decisions, deepening rounds, and given-clause iterations.  The
Marathon alarm provides a final outer bound for code paths whose internal
checks are too sparse.  Timeout returns ``no candidate'' rather than a verdict.

Caches follow the same lifetime discipline.  Parser and static-law data live
for the process.  A superposition run owns its unifier, normalization, term,
and clause caches, which are cleared before another problem.  Marathon could
in principle share more proof state across rows, but the current implementation
does not rely on cross-problem lemmas.  This keeps Solo and Marathon certificate
semantics aligned and limits memory growth under a long manifest.

\section{Countermodel Construction}
\label{sec:false}

For a false implication the solver must find an operation satisfying all
assignments of the hypothesis and at least one assignment falsifying the
goal.  It evaluates candidate families in Python and retains a concrete
witness; the three negative certificate shapes and their common acceptance
boundary are defined in Section~\ref{sec:certarch}.

\subsection{Coefficient matching over structured families}

Represent a magma term by the formal coefficients induced by a candidate
operation.  For
\[
  x \op y = ax+by \pmod n,
\]
each leaf contributes a coefficient obtained by multiplying edge labels along
its path.  Equality of two terms for all assignments is therefore equality of
their coefficient vectors modulo $n$.  This gives an exact, inexpensive test
for whether the hypothesis is an identity of the candidate magma.  The same
representation quickly locates an assignment on which the goal coefficients
differ.  The scan includes composite moduli and its shallow range extends to
the largest carrier admitted by the certificate-instance guard.

Affine operations add a constant term.  Matrix-linear operations replace
scalar coefficients by matrices over small finite fields.  These searches
include two-dimensional vectors over the three-element field and
three-dimensional vectors over the two-element field.  Polynomial candidates
are evaluated on full assignments when the carrier and hypothesis arity make
that feasible.  The deep pass adds a quadratic grid, a two-dimensional family
over the five-element field, and sampled polynomials of degree at most three.
An instance-count guard prevents a mathematically compact operation from
producing a certificate whose exhaustive hypothesis check exceeds the judge
budget.

These families provide both search efficiency and structured certificates.
They find many witnesses faster than a generic table search while preserving
semantic structure in the emitted certificate.  For example, a linear
countermodel can be represented by its arithmetic operation rather than by a
large opaque table.  The judge still evaluates the hypothesis and the failing
goal instance, so coefficient matching remains an untrusted search
optimization rather than a proof rule.

Coefficient evaluation is implemented recursively.  A variable leaf maps to
its basis vector.  At an internal node, the coefficient vector is
$a\vec{u}+b\vec{v}$ for child vectors $\vec{u}$ and $\vec{v}$; affine families
carry one additional constant coordinate.  Thus the universal hypothesis test
does not enumerate assignments.  If the two sides have different vectors,
the candidate is immediately rejected as a model of the hypothesis.  If the
hypothesis vectors agree and the goal vectors differ, a separating assignment
can be sought in the small module.  Full evaluation is still applied before
emission, which protects the certificate path from an error in this symbolic
shortcut.

Vector-linear families use exactly the same recursion with matrices acting on
coordinates.  They are useful because scalar coefficient coincidence may
hold in every small modulus even when non-commuting linear actions distinguish
the goal.  Polynomial operations abandon the linear invariant and instead use
compiled term evaluators over bounded carriers.  The family ordering therefore
progresses from symbolic coefficient equality to increasingly expensive full
evaluation, while presenting a uniform result to the certificate printer: a
carrier, a total binary operation, and a goal witness.

\subsection{Irregular finite models}

Some countermodels have no low-degree algebraic description.  The finite-model
stage treats the Cayley table as a constraint problem.  For a proposed carrier
size, table cells begin unassigned.  Ground instances of the hypothesis watch
the cells on which their two terms depend.  Assigning a cell may simplify a
term, force another cell, or violate an instance.  The search uses a
minimum-remaining-values choice among table positions and abandons partial
tables as soon as a hypothesis instance becomes unequal.

Many identities imply that rows, columns, or both must be permutations.  A law in
which a variable occurs once on one side and appears under the operation in a
particular argument position can expose such a constraint.  The solver infers
the applicable Latin flags, rejects duplicate row or column values, and forces
naked singletons.  For one residual requiring an eight-element quasigroup,
this propagation changed the outcome from no witness within a
one-hundred-twenty-second attempt to a witness in 0.3 seconds.  The production
shallow search remains restricted to carriers four through ten and an eight-second cap; the deep pass revisits the
same carrier interval with a longer cap.

Before emission, the completed table is checked by a separate evaluator over
all hypothesis assignments and a stored goal-falsifying assignment.  The
\lean certificate then reconstructs the table operation and invokes the
judge's finite-decision tactic.  The duplicated checks reduce wasted judge
calls, while only the latter determines acceptance.

The partial-table solver maintains two invariants.  Every assigned cell agrees
with all hypothesis instances that have become ground, and every inferred
Latin constraint agrees with the assigned portion of its row or column.  A
watch list connects a cell to the term instances whose evaluation may advance
when that cell is filled.  Propagation reaches a fixed point before the next
decision.  At a complete table, the first invariant implies that exhaustive
reevaluation of the hypothesis should succeed; the separate final evaluator
checks this implication and also searches the goal assignments.  Search
backtracking may be unsound or incomplete as an algorithm without threatening
an accepted answer, because the table is subsequently evaluated from scratch
and then checked once more in \lean.

Carrier growth is likewise heuristic.  The shallow and deep passes enumerate
the documented carrier interval, but the solver does not infer that absence of
a model there implies truth.  It passes the row to the next proof or
model family.  This behavior implements the paper's no-completeness boundary.

\subsection{Central groupoids and the E168 residual family}

The evaluation-distribution drill identified a family of twelve residuals whose
hypothesis is the central-groupoid law.  Central groupoids have a well-known
combinatorial structure studied by Knuth~\cite{knuthCentral}.  The natural
constructions tried by the structured stages satisfy both the hypothesis and
the relevant E168 conclusions, so they do not separate these equations.
Increasing the generic finder budget did not reliably recover the missing witnesses.

The solver therefore includes one explicit non-natural central groupoid of
order nine.  The table was checked against each relevant hypothesis and goal,
then stored once as a fixed witness family.  When the hypothesis matches up
to variable renaming and a goal is separated by the table, the solver emits
the ordinary finite-table certificate.  This stage is not a classification of
central groupoids and makes no claim about minimal order.  It provides a
specific, auditable witness for a recurring algebraic family: the official
runner accepted the emitted certificates for the twelve drill residuals.

\subsection{Infinite carriers for Austin pairs}

Finite satisfiability can conceal false general implications.  Several pairs
catalogued in the Equational Theories Project are true over the finite magmas
searched by standard methods but false in general.  For recognized Austin-pair
hypotheses, the solver selects a fixed operation on $\Nat$, or its opposite,
and a small tuple falsifying the goal.  The certificate proves the hypothesis
through a fixed sequence of elementary lemmas and closes the negated goal with
the concrete tuple.

Unlike finite decision, this search cannot discharge the universal premise by
enumeration.  Section~\ref{sec:certarch} gives the explicit theorem compiler;
Section~\ref{sec:leanmigration} records the import and compilation costs that
motivated its core-only form.

\subsection{Certificate shapes and the allow-list}
\label{sec:allowlist}

Section~\ref{sec:certarch} specifies the reflected finite shape, the
single-digit limitation of \code{finOpTable}, the transparent namespaced
operation used for larger carriers, and how the dependency inspector treats
local definitions.  The official judge accepted the recorded larger-carrier
certificates in this alternative shape.

\section{Proof-Producing Ordered Superposition}
\label{sec:true}

The true branch operates in unit equational logic.  It assumes one universally
quantified identity and attempts to derive the goal identity.  The engine is
inspired by ordered superposition and unfailing completion as implemented in
first-order provers, but is specialized to one binary symbol, unit clauses,
and the need to reconstruct a small positive \lean proof.

\subsection{Terms, ordering, and inference}

Terms are immutable trees over variables and the binary magma symbol.  A
Knuth--Bendix ordering with unit symbol weights determines maximal equation
sides and orients usable rewrite rules.  Superposition overlaps an oriented
side with a non-variable subterm of another maximal side, computes a most
general unifier, and constructs the resulting equality.  Ordering checks are
repeated after unification, because substitution can change comparisons.
Equations that cannot be oriented remain available for symmetric inferences.

More precisely, all variables in the two premises are renamed apart.  For
equations $l=r$ and $s=t$, a non-variable position $p$ of $s$, and
$\sigma=\mathop{\mathrm{mgu}}(l,s|_p)$ (with an occurs check), the implemented
positive inference has the schematic form
\[
\frac{l=r \qquad s=t}
     {s[r]_p\sigma=t\sigma}.
\tag{SP}
\]
The selected source direction must be KBO-decreasing after instantiation unless
the premise is unorientable; for an unorientable premise the direction is
eligible only when its instantiated right side is not greater.  The selected
target side must be maximal after instantiation, i.e., $t\sigma\not> s\sigma$.
The implementation considers both sides exactly when static KBO cannot orient
the premise.  With unit weights, KBO first checks the variable-occurrence
condition, then total tree size, and finally lexicographic order for the single
binary symbol; fresh goal constants are ordered by their names.

Demodulation is the matching instance of (SP).  If $l\theta=s|_p$, it rewrites
$s$ to $s[r\theta]_p$ only for an oriented rule or when the concrete rewrite is
KBO-decreasing.  Forward demodulation normalizes new equations; a newly active
oriented rule backward-demodulates and retires reducible active equations.
Equations reduced to $q=q$ are deleted.  For a negative goal $u\neq v$, the
corresponding rule is
\[
\frac{l=r \qquad u\neq v}
     {u[r]_p\sigma\neq v\sigma},
\tag{NEG}
\]
with the same source condition and a maximality check on the selected negative
side.  A negative clause closes when its two sides unify.  These rules describe
the implementation rather than asserting a new completeness theorem.

The goal is represented internally as a negative unit clause $u\neq v$ over
fresh constants.  Active equalities can superpose into either side; ordinary
demodulation simplifies the clause.  If its sides unify, the negative clause
is refuted.  This internal negative representation is convenient for search,
but the emitted proof does not ask \lean to validate a refutation calculus.
It instead reconstructs a positive equality chain ending at the original
goal.

Forward demodulation normalizes each generated equality by current rules.
Backward demodulation revisits active equations when a new rule can simplify
them.  Symmetric alpha-normalized keys remove variants.  A given-clause loop
interleaves a weight heap with an age queue; later deepening rounds alter the
variable penalty and age ratio so that a proof blocked by one selection policy
is not permanently starved.

\subsection{The tautology-deletion failure}

An early version of the prover did not resolve one public residual,
\code{hard3\_0314}.  The official-runner attempt found no proof after 810 seconds
even though an external derivation suggested that the necessary clauses were within the
configured size bound.  Instrumentation indicated that throughput was not the
primary cause.  When demodulation transformed an equation into $s=s$, the
step constructor returned no proof node, and the caller retained the partially
rewritten input equation.  Backward simplification had the same behavior.

Consequently, many small instances of the hypothesis remained live.  Their
weights made them attractive to given-clause selection, so they repeatedly
displaced clauses on the useful derivation.  The correct simplification rule
is deletion: a demodulated tautology contributes nothing to saturation.  The
fix retires the active clause before testing the rewritten result and drops a
new clause whose sides become equal.

With that change, the prover found the residual proof in 0.6 seconds in its
direct measurement, and the official-runner path produced an accepted
certificate in 14.6 seconds of wall-clock time.  These measurements illustrate
the coupling between simplification and selection rather than establishing a
general speed ratio.  Retaining a tautology as a live clause changed the
effective search strategy enough to prevent the final public proof from being
found within the earlier attempt.

\subsection{Indexing and allocation control}

Demodulation candidates are retrieved through a discrimination tree keyed by
a term's linear skeleton.  The index returns rules whose left sides may match a
subterm; matching and the ordering check filter false positives.  Normal forms
are cached and versioned by the rewrite-rule set.  Term sizes, variable
multisets, preorder traversals, and eligible positions are cached on immutable
equations.

The unifier returns a triangular substitution rather than eagerly applying it
to every binding.  A memoised substitution routine resolves chains while
sharing unchanged subtrees.  Before constructing a superposition result, the
engine estimates instantiated size from static term sizes and per-variable
counts.  Overlaps that cannot fit the current cap are rejected before
allocating their terms.  Unifier results are cached by the renamed overlap
pair.  These changes preserve the inference trace at fixed bounds while
reducing the work per surviving clause.

The passive set stores recipes---source equation, target equation, direction,
and overlap position---rather than full instantiated proof objects.  The
selected recipe is materialized when it becomes given.  Under the hosted
memory limit, this representation permits a large passive frontier without
duplicating every term and ancestry chain.  Caches and negative-goal storage
are bounded and cleared between prover invocations.

\subsection{Anytime deepening}

A fixed size cap is sensitive to the selected bound: saturation at a small cap
can terminate with time remaining even when one slightly larger intermediate would finish the proof.
The deep pass therefore runs a schedule of increasing term-size and variable
caps.  If a round saturates, its unused time moves forward.  If the final
configured round saturates, caps continue to grow within the overall deadline.
The last strategies reduce the variable penalty and increase age pressure,
approximating a different prover configuration without maintaining another
implementation.

The procedure remains budget-bounded and machine-load sensitive.  It is not a
new completeness result for ordered superposition, nor a proof that the chosen
ordering is fair under every cap.  The evaluation in
Section~\ref{sec:evaluation} measures this implementation on fixed corpora.

\subsection{Replay into the \lean kernel}

Section~\ref{sec:certarch} gives the complete compiler: reachable-proof
slicing, substitution and context reconstruction, topological emission,
positive-chain conversion, robust re-emission, the replay invariant, and a
worked Lean term.  The search engine described in this section supplies its
recipe DAG; it never asks Lean to trust the ordering or refutation calculus.

\section{Engineering for the Hosted Environment}
\label{sec:hosted}

\subsection{Toolchain migration}
\label{sec:leanmigration}

Section~\ref{sec:certarch} gives the final core-and-judge-module dependency
policy.  Historically, the public harness used \lean 4.30.0-rc2, while the
announced hosted verifier used \lean 4.32.0 with the corresponding Mathlib
release.  Most emitted certificates already followed that policy.  The
exception was the infinite-carrier Austin family, whose tactic import made
compilation sensitive to the larger library environment.  In a compatibility
run, one such certificate required about 330 seconds of compilation.

Rather than increasing the solver-side timeout, we rewrote the Austin
certificate in positive form, replaced the broad tactic dependency with
explicit core lemmas, and re-emitted the witness.  Its \lean 4.32 compilation
then completed within 3.1 seconds.  A stratified corpus of 120 accepted certificates, covering every emission
family and all fourteen former v1 residuals, compiled 119 of 120 under the
exact announced toolchain, the single non-pass being the pre-rewrite witness
itself at an external 120-second phase limit that the judge's own 300-second
configuration clears.  That corpus characterizes the migration-era certificate
set rather than the frozen submission's current output, which for the same
problem imports only the judge module and is accepted in under five seconds.
The judge support modules compiled unchanged; the run is archived as an
operator attestation (\code{results/lean432\_corpus/}).  When the organizers subsequently upgraded
the judge to \lean 4.33.1 after the kernel-soundness review, the frozen
artifact was revalidated in full under the upgraded judge --- official harness,
all six public sets, the distribution drill, and the Marathon manifest --- with
the resulting ledgers committed and hash-bound
(\code{results/lean4331\_revalidation/}).  Neither measurement substitutes for
hosted scoring.

\subsection{Input encoding}
\label{sec:encoding}

The official judge normalizes both operation spellings when it builds the
\lean problem.  The runner, however, passes the original problem text to
the contestant process.  An earlier solver normalized only along one internal
path.  When the Stage 1 evaluation splits were supplied in their published
asterisk encoding, the solver crashed before certificate search and recorded
zero accepted rows out of 800.

The frozen solver normalizes both equations at both the Solo and Marathon intakes.
The same drill then recorded 800 accepted certificates with full agreement
against the published answers.  Thus, verified output does not address failures
in an unverified parser; normalization must occur at the protocol boundary,
before any dispatch or caching.

\subsection{Resource and filesystem constraints}

The hosted-style sandbox has a read-only submission mount, a bounded process
count, bounded memory, no direct network, and a small writable temporary
filesystem.  The solver consequently uses the Python standard library, keeps
all static data in the single source file, does not spawn external provers, and
does not assume repository access.  Bounded passive queues and caches keep the
superposition engine below the sandbox memory ceiling observed in long false
saturations.  The process-level timer also prevents a deeply nested legacy
loop from overrunning its Marathon slice.

Marathon's output file is the only durable mutable state supplied by the
runner, so per-row flush and synchronization are required for correctness; a
managed workspace whose default artifact directory was unwritable was resolved
by redirecting judge artifacts to temporary storage, leaving judge sources
untouched.

\section{Evaluation}
\label{sec:evaluation}

\subsection{Method and provenance discipline}

Unless marked hosted, measurements in this section are local runs of the
official runner and \lean judge.  The main regression used
official repository revision \texttt{2848228ff490...}; the harness completed
without errors on the operator host.  Table~\ref{tab:identity} is the sole
artifact identity used by every result below.  Each result file is bound by the
repository provenance manifest; short hashes in later tables are prefixes of
the full SHA-256 values in \code{PROVENANCE.json}.

\begin{table}[t]
\caption{Canonical artifact identity.  The freeze timestamp precedes the
official stress-set release dated 2026-08-28.}
\label{tab:identity}
\small
\begin{tabularx}{\columnwidth}{@{}lY@{}}
\toprule
Field & Canonical value \\
\midrule
Solver & v2.5, 189,504 bytes \\
SHA-256 & \scriptsize\texttt{f2392533c9f4c03b292be80bc6d12e98}\newline
  \scriptsize\texttt{e5254cc4861d1cc4b227957ad5ed89b4} \\
Freeze & \scriptsize\texttt{7b636bc3c3e8f4eed80fea1809ba940af3c466f0}, 2026-08-26 21:31:32 UTC \\
Local judge & \scriptsize\texttt{2848228ff490422442878fd6f5abaf4cfa95257d}; Lean 4.30.0-rc2, Mathlib \texttt{896cc56a} \\
Hosted recheck & \scriptsize\texttt{13648682a5553717ea91b86513ed140b39160cf5}; Lean 4.33.1, Mathlib \texttt{0df444a3} \\
\bottomrule
\end{tabularx}
\end{table}

The reported wall-clock time is the sum of per-row runner measurements, not
elapsed parallel makespan and not a hardware-independent cost model.  Solver stages
with deadlines may take different paths under heavier load.  ``LLM'' counts
language-model calls attributed by the runner.  All reported accepted results
are produced deterministically, with zero language-model and zero
external-prover calls.

The Solo path retains a fallback that calls a language model only after the
certificate stages; Marathon disables it, and any Solo proposal would still
pass through the same printer and judge.  No archived accepted row invoked
it.\footnote{A separate historical measurement of that configuration, before
its optimization, is recorded in \code{PROVENANCE.json}: the organizer-pinned
model settled none of six residuals and varied between first attempts under
fixed temperature and seed.  It diagnoses one configuration and supports no
general claim.}

\subsection{Six public sets}

Table~\ref{tab:public} reports the final six-set regression.  All rows were
accepted, every accepted row had one judge call, and no row used the fallback
model.  The sum of the set sizes is 1,889.  The earlier frozen v1 solver
accepted 1,875 of those rows; v2 adds the structured false stages and G3 proof
engine described above.  We report v1 only as development context, not as a
comparison against another entrant.

\begin{table*}[t]
\caption{Local official-runner regression for solver v2.5.  ``Wall'' is the
sum of per-row wall-clock seconds.  Each row names its immutable ledger and
SHA-256 prefix.}
\label{tab:public}
\small
\setlength{\tabcolsep}{4pt}
\begin{tabularx}{\textwidth}{@{}lrrrrY@{}}
\toprule
Set & Accepted & LLM & Judge & Wall (s) & Ledger provenance \\
\midrule
\code{sample\_20}  & 20/20 & 0 & 20 & 67.53 & \prov{results/v2\_sample\_20\_official\_2848228.json}{b17eaff452d2...} \\
\code{sample\_200}  & 200/200 & 0 & 200 & 964.12 & \prov{results/v2\_sample\_200\_official\_2848228.json}{dcf830de6199...} \\
\code{hard1}  & 69/69 & 0 & 69 & 559.81 & \prov{results/v2\_hard1\_official\_2848228.json}{7724566d1ceb...} \\
\code{hard2}  & 200/200 & 0 & 200 & 1250.67 & \prov{results/v2\_hard2\_official\_2848228.json}{3e68548945d4...} \\
\code{hard3}  & 400/400 & 0 & 400 & 2857.38 & \prov{results/v2\_hard3\_official\_2848228.json}{c9f98750e018...} \\
\code{normal}  & 1000/1000 & 0 & 1000 & 4128.3 & \prov{results/v2\_normal\_official\_2848228.json}{27c756482d15...} \\
\bottomrule
\end{tabularx}
\end{table*}

The totals establish regression coverage of these public artifacts only.  The
private evaluation set is separate, and public-set tuning is possible.  We do
not treat Table~\ref{tab:public} as a hidden-set estimate.

\subsection{Certificate and checking cost}

Table~\ref{tab:certcost} first counts the accepted Lean artifacts in the
principal evaluation and then aggregates UTF-8 source bytes and accepted
judge-event time from the bound public and holdout ledgers.  ``Check'' includes
judge orchestration around Lean compilation and kernel checking; it is not
prover search time.  Size and time distributions are separated by polarity
because reflected countermodels and equality-chain proofs have different
shapes.

\begin{table*}[t]
\caption{Accepted Lean artifacts and certificate costs.  Panel (a) binds each
empirical corpus to its source; the total is their arithmetic sum.  In panel
(b), entries are median / 95th percentile / maximum.  Public and holdout
statistics derive from \code{results/paper\_certificate\_metrics.json},
SHA-256 \texttt{56f233426fd63009...}; its embedded source hashes include the
holdout ledger \texttt{c741609293d136c6...}.}
\label{tab:certcost}
\scriptsize
\begin{tabularx}{\textwidth}{@{}p{1.28in}rrrrrY@{}}
\toprule
\multicolumn{7}{@{}l}{\textbf{(a) Accepted Lean artifact production}} \\
Corpus & \shortstack{Positive Lean\\proofs} &
\shortstack{Negative Lean\\witnesses} & Total & LLM calls &
\shortstack{External\\prover} & Ledger provenance \\
\midrule
Six public sets & 929 & 960 & 1,889 & 0 & 0 &
\prov{results/paper\_certificate\_metrics.json}{56f233426fd63009...} \\
Temporal holdout (official stress set) & 100 & 100 & 200 & 0 & 0 &
\prov{results/stage2\_stress\_test\_results.json}{c741609293d136c6...} \\
Published distribution drills & 400 & 400 & 800 & 0 & 0 &
\prov{results/official\_distribution\_drill/*.json}{6647e3a13fc0... / 2025ca90a892... / b931bbf33007... / d7ef36ee24f7...} \\
\textbf{Total} & \textbf{1,429} & \textbf{1,460} & \textbf{2,889} &
\textbf{0} & \textbf{0} & Derived from the three hash-bound rows above \\
\bottomrule
\end{tabularx}
\vspace{5pt}

\begin{tabular}{@{}llrrr@{}}
\toprule
\multicolumn{5}{@{}l}{\textbf{(b) Certificate source and checking cost}} \\
Corpus & Verdict ($n$) & Certificate bytes & Check time (s) & End-to-end wall (s) \\
\midrule
Public & false (960) & 306 / 498 / 6,885 & 2.20 / 8.85 / 41.74 & 2.30 / 13.43 / 41.95 \\
Public & true (929) & 494 / 1,392 / 10,766 & 1.43 / 6.98 / 23.21 & 4.26 / 15.43 / 35.29 \\
Temporal holdout & false (100) & 323 / 551 / 551 & 2.64 / 14.89 / 17.92 & 2.77 / 15.69 / 18.04 \\
Temporal holdout & true (100) & 500 / 1,406 / 12,485 & 1.72 / 6.13 / 10.04 & 5.91 / 12.12 / 66.49 \\
\bottomrule
\end{tabular}
\end{table*}

Peak memory was not recorded by the official runner and is therefore not
reported.  The bounded recipe queues described in Section~\ref{sec:true} are a
design response to the hosted limit, not a substitute for measured resident
set size.

\subsection{Component-level interpretation}

The superseded v1 baseline left fourteen public residuals: eleven true
implications without a proof inside its budgets and three false
implications without a model in its available families and carriers.  The v2
development was organized around those failure classes.  G3 plus the
tautology-deletion correction supplies certificates for the true residuals;
the large-carrier, Latin-propagating, structured, and infinite-carrier paths
supply witnesses for the false residuals.  Replaying the residual set alone
records fourteen accepted certificates, whereas the six complete sets provide
a regression check that stage ordering and printers preserve the previously
accepted cases.

This history is diagnostic, not a controlled ablation study.  Several changes
were introduced together, and a row may now be solvable by more than one stage.
The paper therefore attributes mechanisms only where the engineering notes contain
a direct replay, such as \code{hard3\_0314} for tautology deletion and
\code{hard1\_0062} for Latin propagation.  It does not assign an aggregate
percentage gain to individual components.

The order of the cascade also affects the observed wall totals.  A true row
may incur the shallow false-family cost before receiving a quick proof; a false
row may incur direct proof checks before finite-model search.  These costs are
intentional because the early stages reduce resource use across both branches:
inexpensive structured countermodels avoid proof saturation, while direct proof
patterns avoid table search.  The six-set totals measure the integrated policy,
not isolated engine speed.

\subsection{Distribution drill, Marathon, and hosted playground}

The official-distribution drill uses the four published Stage 1 evaluation
splits that were named as Stage 2 scoring categories.  Each split contains
published ground truth and uses the asterisk input encoding.  Table
\ref{tab:drill} reports full local judge acceptance and full agreement with
those labels.  Stage 2 does not reuse these problems, so the drill measures
format and distribution readiness rather than evaluation transfer.

The canonical Marathon measurement used the official batch runner and its
\code{normal\_100} manifest.  All rows were accepted without tokens.  The
hosted-playground row is different in kind: it was run through the official
playground interface on the hosted judge and is recorded as an external
operator attestation, because the playground emits no repository-committable
ledger.  Across the four evaluation categories the attested runs record
acceptance of every attempted problem with no rejected or errored rows and no
language-model calls.  It is explicitly not an official leaderboard score.

\begin{table*}[t]
\caption{Additional measurements.  Local drill and Marathon rows name their
ledger files and SHA-256 prefixes.  The hosted row is an operator
attestation recorded in the \code{hosted\_playground\_measurement} entry of
\code{PROVENANCE.json}; the playground interface did not emit a repository
ledger file.}
\label{tab:drill}
\small
\setlength{\tabcolsep}{4pt}
\begin{tabularx}{\textwidth}{@{}lrrrY@{}}
\toprule
Measurement & Accepted & LLM/tokens & Wall statistic & Ledger provenance \\
\midrule
\code{evaluation\_order5} & 200/200 & 0 calls & sum 2047.27 s & \prov{results/official\_distribution\_drill/evaluation\_order5\_results.json}{6647e3a13fc0...} \\
\code{evaluation\_extra\_hard} & 200/200 & 0 calls & sum 17307.45 s & \prov{results/official\_distribution\_drill/evaluation\_extra\_hard\_results.json}{2025ca90a892...} \\
\code{evaluation\_hard} & 200/200 & 0 calls & sum 2228.57 s & \prov{results/official\_distribution\_drill/evaluation\_hard\_results.json}{b931bbf33007...} \\
\code{evaluation\_normal} & 200/200 & 0 calls & sum 1019.91 s & \prov{results/official\_distribution\_drill/evaluation\_normal\_results.json}{d7ef36ee24f7...} \\
Marathon \code{normal\_100} & 100/100 & 0 tokens & full default budget & \prov{results/marathon/normal\_100\_summary.json}{d422d4ad6703...} \\
Hosted playground (\code{evaluation\_normal}) & 200/200 & 0 calls & mean 4.89 s/problem & \prov{PROVENANCE.json: hosted\_playground\_measurement}{artifact 7916cbc2...} \\
\bottomrule
\end{tabularx}
\end{table*}

The hosted record further reports means by verdict: 2.69 seconds for
false rows and 7.08 seconds for true rows, with maxima 6.5 and 23.8 seconds,
respectively.  Those values include hosted orchestration and certificate
checking as exposed by the playground, and should not be compared directly to
the local prover-only benchmark below.

\subsection{Temporal holdout and research frontier}
\label{sec:holdout}

The solver artifact in Table~\ref{tab:identity} was frozen at 21:31:32 UTC on
2026-08-26.  The organizers released the official Stage 2 stress
set on 2026-08-28, after that freeze.  The set was
announced as mirroring the final leaderboard composition: four categories of
fifty rows, each containing twenty-five true and twenty-five false
implications.  The byte-identical solver then produced 200/200 accepted
certificates with zero language-model calls in 22 minutes 25 seconds.  Its
per-category verdict counts match the announced composition, although the set
does not publish per-item reference labels.  Its ledger has SHA-256 prefix
\texttt{c741609293d136c6}; the full digest is bound in \code{PROVENANCE.json}.

The temporal ordering, immutable solver hash, and absence of post-release
parameter changes make this a holdout with respect to solver development.  It
is not a random sample from all magma implications, and its announced
competition distribution may resemble the public suites, but unlike those
suites no row could have driven a solver change.

The simultaneously released, unscored order-five research set gives the
opposite result.  It contains 100 research-tier problems with an
order-five Austin law on at least one side, partly unknown ground truth, and no
finite countermodel for the relevant false directions.  The artifact
certified 0/100 within budget: 73 searches exhausted the deep budget and 27
candidate certificates were rejected by the kernel.  This negative ledger
(SHA-256 \texttt{86fc2764fdb9...}) prevents the holdout result from being read
as broad completeness.

\section{Honest Boundary and Reproducibility}
\label{sec:repro}

\subsection{Scope of the empirical claims}

The scientific claims do not depend on a leaderboard score or rank.  Public
and distribution-drill rows measure regression, the hosted playground measures
toolchain compatibility, and Section~\ref{sec:holdout} supplies the temporal
holdout; none of them is a completeness theorem or a cross-system comparison.
The coverage is bounded by construction: structured countermodels cover chosen
families, finite-model search has carrier and time bounds, and superposition
deepens only until its deadline.  Which problems the deep stages reach can
vary with machine load; the judge's acceptance of a fixed certificate under a
pinned toolchain does not.

\subsection{Threats to validity}

The public suites influenced solver development.  The fourteen v1 residuals
motivated particular prover and countermodel changes, and the E168 drill
motivated the fixed central groupoid.  Full regression shows that those
changes integrate without losing earlier rows; it does not measure performance
on an independently sampled corpus.  The distribution drill is broader but is
still published Stage 1 data, and its ground truth may share construction
patterns with the public suites.

Wall-clock totals combine Python search, process overhead, and \lean checking
under one operator-host environment.  They are appropriate for detecting large
regressions within that setup but should not be read as portable runtimes.
Parallel worker load affects time-budgeted deepening, and filesystem caching
affects \lean startup.  The hosted-playground mean times partially address this
environment gap, but playground selection and daily interaction are not the
same protocol as final evaluation.

Finally, kernel checking validates each emitted theorem relative to the
judge's imported environment and allowed axioms.  It does not validate that
the benchmark labels, competition problem generator, or provenance narrative
are correct.  We mitigate the last concern with hashes and exact commands;
benchmark construction and judge implementation remain external assumptions.

\section{Related Work}
\label{sec:rw}

\paragraph{Equational theories.}
The Equational Theories Project builds a collaborative implication graph for
single magma identities and supplies formal proofs, finite models, and curated
problem data~\cite{etp}.  The SAIR task derives its laws and implication
setting from this project but imposes a distinct online solver and certificate
interface.  Cazares analyzes Stage 1 of the challenge and its prediction
setting~\cite{cazaresStage1}.  Stage 2 changes the observable output from a
label to a judge-accepted proof or countermodel.

Knuth's study of central groupoids underlies the E168 witness
family~\cite{knuthCentral}; the use here is narrower than that structural
theory, embedding one checked non-natural finite model because natural
constructions separated none of the goals encountered.  The Austin models
likewise draw on constructions represented in the Equational Theories Project
rather than proposing a general finite-model theorem.

\paragraph{Proof-producing superposition and reconstruction.}
Duper is the closest direct precedent, implementing proof-producing
superposition as a Lean tactic for dependent type theory~\cite{duper}; it is
broader in logic and tighter in proof-assistant integration, and sets the
positive baseline here.  The complementary point studied here couples a
narrower equality compiler to finite and infinite negative witness compilers
under one dependency policy and a single-file scheduler.

Lean-Auto and Lean on Vampire translate Lean goals to external provers and
reconstruct their answers~\cite{leanAuto,leanVampire}, answering the natural
alternative ``call a mature ATP, then replay.''  That architecture supplies a
stronger general prover; the deployment contract here admits a single Python
file, and specialization buys proof search the budget to share parsing,
scheduling, and certificate dispatch with the countermodel branch.

Isabelle's Sledgehammer architecture delegates search to external provers and
uses tactics such as Metis for reconstruction~\cite{sledgehammer,metis}.
The \textsc{slam} tactic instead targets higher-order logic directly with
$\lambda$-superposition and can itself serve as a Sledgehammer backend
\cite{slam}.  Our positive replay shares the same skeptical principle, but it
emits first-order unit-equality chains rather than translating higher-order
proofs.  TSTP provides a general interchange vocabulary for ATP problems and
derivations~\cite{tstp}; our recipes are smaller and domain-specific, but this
choice sacrifices interoperability with mature ATP tooling.

\paragraph{Formalized calculi and model finding.}
Formalizations of superposition study the soundness and completeness of the
calculus itself; recent Isabelle work extends such a framework with sorts
\cite{sortedSuperposition}.  Our Python calculus is not formally verified.
Instead, each successful derivation is compiled away and checked extensionally
as a Lean equality proof.  This reduces the trusted base for each answer but
does not establish completeness or correctness of the search algorithm.

Vampire~\cite{vampire}, E~\cite{eprover}, and Prover9/Mace4
\cite{prover9} implement mature saturation and finite-model architectures;
Paradox develops efficient MACE-style finite model finding
\cite{paradox}.  Our bounded finder is less general.  Its system contribution
is that the discovered table is not accepted as a Python-side Boolean: it is
embedded in a negative Lean witness and exhaustively checked under the same
contract as a positive proof.  Structured and infinite families cover cases
where generic small-table enumeration is either wasteful or inapplicable.

\paragraph{Lean integration.}
Section~\ref{sec:certarch} collects the Lean-specific consequences of equality
elimination, finite reflection, dependency inspection, and toolchain startup
\cite{lean4,mathlib}.  Together they favor small, locally typed certificates
over replaying a global search trace.  A reusable length-independent
finite-operation certificate API would remove the remaining
proof-assistant-specific printer branch.

\section{Conclusion}

The SAIR EQT2 Stage 2 contract rewards an answer only when the same artifact
also supplies a machine-checked reason.  The solver described here meets that
contract with one fail-closed cascade for both polarities: structured and
irregular countermodels, infinite witnesses, and ordered unit superposition
all terminate in judge-checked Lean terms.  This dual-polarity integration,
the restricted single-file deployment, and the evidence discipline distinguish
the system from proof-producing superposition alone.  Public regressions show
integration coverage; the 200/200 temporal holdout tests the byte-identical
artifact after freeze, while the research-tier failure marks a concrete
frontier.
The broader lesson is that aggressive untrusted search remains auditable when
each stage preserves enough structure to emit a small positive proof or an
explicit negative model, and when empirical claims are bound to the exact
bytes that produced them.

\ifplainarticle
  \bibliographystyle{plainnat}
\else
  \bibliographystyle{ACM-Reference-Format}
\fi
\ifanonsubmission\else
\section*{Author contributions (CRediT)}
Haobo Ma: conceptualization, methodology, software (solver architecture and the
systems foundation), investigation, validation.
Wenlin Zhang: conceptualization, methodology, software, investigation,
validation, data curation, project administration, writing (original draft).
Manuel Israel C\'azares: validation (solver freeze verification --- hash, file
size, and single-file contract compliance; evidence-ledger auditing across the
reviewed freezes); methodology (review of experimental design and the
claims--evidence boundary); writing (review and editing --- full manuscript
review and critical review of the presentation and claims).
\fi

\bibliography{references}

\appendix

\section{Reproducibility Materials}
\label{app:repro}

\subsection{Hash-bound evidence}

\code{PROVENANCE.json} binds the frozen solver's SHA-256 and byte count, the
official repository and configuration revisions, each result-ledger hash, the
Marathon outputs, the compatibility corpus, and the hosted-playground record.
The human-readable claims ledger states the maximum set of supported claims
and separately lists prohibited inferences.  The repository script
\code{scripts/check\_freeze.py} recomputes the solver and ledger hashes, checks
accepted-row metadata, verifies the one-file submission layout, and fails
closed on drift.

This structure prevents stale measurements from being attributed to a
different solver.  Every table uses the v2.5 identity in
Table~\ref{tab:identity}; older provenance entries are outside the paper's
quantitative claim surface.

\subsection{Replay procedure}

Reproduction requires this artifact and the official Stage 2 judge at the
pinned revision.  The artifact-side integrity check is
\begin{verbatim}
python3 scripts/check_freeze.py
\end{verbatim}
The pinned Solo or Marathon runner then consumes a submission directory
containing only the frozen \code{solver.py}.
The anonymous supplement includes the exact pinned revisions, manifests,
ledgers, and replay commands.  Recomputed output must be stored separately
rather than overwriting the archived evidence.

The \lean 4.32 compatibility corpus was run in an external directory and is
described in the provenance record; it is not presently packaged as a
one-command repository test.  This is a reproducibility limitation.  The
official harness, freeze checker, ledger files, and public-runner commands are
available locally; hosted playground interaction and any future leaderboard
submission require organizer infrastructure.


\end{document}
